% GENERATED FILE. Edit canonical lessons, then run npm run generate:books.
% canonical-source-hash: fnv1a64:99f1068348bdd549
% canonical-lessons: ML-C510-mutalmannu, ML-C510-alippalam, ML-C510-mutikkettuka, ML-C510-tanukkuka, ML-C510-irampuka

\chapter{Fog and Hail}
\label{ch:ml-fog-and-hail}
\hlchaptermodality{510}

\begin{chapteropening}
\textbf{By the end of this chapter:} \emph{I can say fog, hail, to cloud over, to grow cold and to rumble.}

Complete the last lesson of chapter 510: I can say fog, hail, to cloud over, to grow cold and to rumble.
\end{chapteropening}

\section[\texorpdfstring{\ml{മൂടൽമഞ്ഞ്} (mūṭalmaññŭ)}{മൂടൽമഞ്ഞ് (mūṭalmaññŭ)}]{\ml{മൂടൽമഞ്ഞ്} (mūṭalmaññŭ) — fog}

\label{lesson:ML-C510-mutalmannu}



\begin{quote}
Before the new one: say the Malayalam for a chill, then the Malayalam for to get hot.
\end{quote}

\subsection*{You'll want to know: \ml{മൂടൽമഞ്ഞ്}}
\textbf{\ml{മൂടൽമഞ്ഞ്}} — \emph{mūṭalmaññŭ} — \textquotedblleft{}fog\textquotedblright{}.

\begin{quote}
\ml{മൂടൽമഞ്ഞ്} \ml{കാരണം} \ml{വഴി} \ml{കാണാൻ} \ml{കഴിഞ്ഞില്ല}.
\end{quote}

\begin{quote}
\emph{mūṭalmaññŭ kāraṇaṁ vaḻi kāṇān kaḻiññilla.}
\end{quote}

\begin{quote}
Because of the fog we could not make out the road.
\end{quote}

\ml{മൂടൽ}, a covering, from \ml{മൂടുക}, plus \ml{മഞ്ഞ്}, mist.

\begin{cousinweb}
\textbf{\ml{മൂടൽമഞ്ഞ്}} (\emph{mūṭalmaññŭ}) is a native Dravidian compound of \ml{മൂടൽ}, a covering, from \ml{മൂടുക} (DEDR 5034), and \ml{മഞ്ഞ്}, mist (DEDR 4641).
\end{cousinweb}

\begin{grammarlens}[title={What you've built}]
One more word for getting around.
\end{grammarlens}

\subsection*{Your turn}
\begin{itemize}
\raggedright
  \item \emph{Say it:} \emph{mūṭalmaññŭ}
  \item \emph{Say it:} \emph{mūṭalmaññŭ}, once more
  \item \emph{Say it:} \emph{mūṭalmaññŭ kāraṇaṁ vaḻi kāṇān kaḻiññilla.} — Because of the fog we could not make out the road.
  \item \emph{Recall:} say the Malayalam for a chill, then the Malayalam for to get hot, then say \emph{mūṭalmaññŭ} again
\end{itemize}

\begin{tcolorbox}[breakable,colback=teal!4,colframe=teal!35!black,title={Before you move on}]
What does \ml{മൂടൽമഞ്ഞ്} mean? (\textquotedblleft{}Fog\textquotedblright{}.) Say it once more.
\end{tcolorbox}
\section[\texorpdfstring{\ml{ആലിപ്പഴം} (ālippaḻaṁ)}{ആലിപ്പഴം (ālippaḻaṁ)}]{\ml{ആലിപ്പഴം} (ālippaḻaṁ) — hail; a hailstone}

\label{lesson:ML-C510-alippalam}



\begin{quote}
Before the new one: say the Malayalam for to get hot, then the Malayalam for fog.
\end{quote}

\subsection*{You'll want to know: \ml{ആലിപ്പഴം}}
\textbf{\ml{ആലിപ്പഴം}} — \emph{ālippaḻaṁ} — \textquotedblleft{}hail; a hailstone\textquotedblright{}.

\begin{quote}
\ml{ഇന്നലെ} \ml{വൈകുന്നേരം} \ml{ആലിപ്പഴം} \ml{വീണു}.
\end{quote}

\begin{quote}
\emph{innale vaikunnēraṁ ālippaḻaṁ vīṇu.}
\end{quote}

\begin{quote}
Yesterday evening there was hail.
\end{quote}

Literally hail fruit: \ml{ആലി}, hail, plus \ml{പഴം}, a fruit.

\begin{cousinweb}
\textbf{\ml{ആലിപ്പഴം}} (\emph{ālippaḻaṁ}) is native Dravidian; DEDR lists \ml{ആലി} and \ml{ആലിപ്പഴം} as hail (DEDR 384), with \ml{പഴം}, a fruit.
\end{cousinweb}

\begin{grammarlens}[title={What you've built}]
One more word for getting around.
\end{grammarlens}

\subsection*{Your turn}
\begin{itemize}
\raggedright
  \item \emph{Say it:} \emph{ālippaḻaṁ}
  \item \emph{Say it:} \emph{ālippaḻaṁ}, once more
  \item \emph{Say it:} \emph{innale vaikunnēraṁ ālippaḻaṁ vīṇu.} — Yesterday evening there was hail.
  \item \emph{Recall:} say the Malayalam for to get hot, then the Malayalam for fog, then say \emph{ālippaḻaṁ} again
\end{itemize}

\begin{tcolorbox}[breakable,colback=teal!4,colframe=teal!35!black,title={Before you move on}]
What does \ml{ആലിപ്പഴം} mean? (\textquotedblleft{}Hail; a hailstone\textquotedblright{}.) Say it once more.
\end{tcolorbox}
\section[\texorpdfstring{\ml{മൂടിക്കെട്ടുക} (mūṭikkeṭṭuka)}{മൂടിക്കെട്ടുക (mūṭikkeṭṭuka)}]{\ml{മൂടിക്കെട്ടുക} (mūṭikkeṭṭuka) — to cloud over}

\label{lesson:ML-C510-mutikkettuka}



\begin{quote}
Before the new one: say the Malayalam for fog, then the Malayalam for hail.
\end{quote}

\subsection*{You'll want to know: \ml{മൂടിക്കെട്ടുക}}
\textbf{\ml{മൂടിക്കെട്ടുക}} — \emph{mūṭikkeṭṭuka} — \textquotedblleft{}to cloud over\textquotedblright{}.

\begin{quote}
\ml{ഉച്ചയ്ക്ക്} \ml{ആകാശം} \ml{മൂടിക്കെട്ടി}.
\end{quote}

\begin{quote}
\emph{uccaykkŭ ākāśaṁ mūṭikkeṭṭi.}
\end{quote}

\begin{quote}
At noon the sky clouded over.
\end{quote}

Past: \ml{മൂടിക്കെട്ടി}, \emph{mūṭikkeṭṭi}.

The sky is the subject: \ml{ആകാശം} \ml{മൂടിക്കെട്ടി}. Literally covered and tied up.

\begin{cousinweb}
\textbf{\ml{മൂടിക്കെട്ടുക}} (\emph{mūṭikkeṭṭuka}) is a native Dravidian compound of \ml{മൂടുക}, to cover (DEDR 5034), and \ml{കെട്ടുക}, to tie (DEDR 1147).
\end{cousinweb}

\begin{grammarlens}[title={What you've built}]
One more word for getting around.
\end{grammarlens}

\subsection*{Your turn}
\begin{itemize}
\raggedright
  \item \emph{Say it:} \emph{mūṭikkeṭṭuka}
  \item \emph{Say it:} \emph{mūṭikkeṭṭuka}, once more
  \item \emph{Say it:} \emph{uccaykkŭ ākāśaṁ mūṭikkeṭṭi.} — At noon the sky clouded over.
  \item \emph{Recall:} say the Malayalam for fog, then the Malayalam for hail, then say \emph{mūṭikkeṭṭuka} again
\end{itemize}

\begin{tcolorbox}[breakable,colback=teal!4,colframe=teal!35!black,title={Before you move on}]
What does \ml{മൂടിക്കെട്ടുക} mean? (\textquotedblleft{}To cloud over\textquotedblright{}.) Say it once more.
\end{tcolorbox}
\section[\texorpdfstring{\ml{തണുക്കുക} (taṇukkuka)}{തണുക്കുക (taṇukkuka)}]{\ml{തണുക്കുക} (taṇukkuka) — to grow cold; to cool down}

\label{lesson:ML-C510-tanukkuka}



\begin{quote}
Before the new one: say the Malayalam for hail, then the Malayalam for to cloud over.
\end{quote}

\subsection*{You'll want to know: \ml{തണുക്കുക}}
\textbf{\ml{തണുക്കുക}} — \emph{taṇukkuka} — \textquotedblleft{}to grow cold; to cool down\textquotedblright{}.

\begin{quote}
\ml{സന്ധ്യയായപ്പോൾ} \ml{കാറ്റ്} \ml{തണുത്തു}.
\end{quote}

\begin{quote}
\emph{sandhyayāyappōḷ kāṟṟŭ taṇuttu.}
\end{quote}

\begin{quote}
When dusk came, the wind grew cold.
\end{quote}

Past: \ml{തണുത്തു}, \emph{taṇuttu}.

For weather, food and drink: \ml{ചായ} \ml{തണുത്തു}, the tea went cold.

\begin{cousinweb}
\textbf{\ml{തണുക്കുക}} (\emph{taṇukkuka}) is native Dravidian, \textquotedblleft{}to grow cold, to be cooled\textquotedblright{} (DEDR 3045).
\end{cousinweb}

\begin{grammarlens}[title={What you've built}]
One more word for getting around.
\end{grammarlens}

\subsection*{Your turn}
\begin{itemize}
\raggedright
  \item \emph{Say it:} \emph{taṇukkuka}
  \item \emph{Say it:} \emph{taṇukkuka}, once more
  \item \emph{Say it:} \emph{sandhyayāyappōḷ kāṟṟŭ taṇuttu.} — When dusk came, the wind grew cold.
  \item \emph{Recall:} say the Malayalam for hail, then the Malayalam for to cloud over, then say \emph{taṇukkuka} again
\end{itemize}

\begin{tcolorbox}[breakable,colback=teal!4,colframe=teal!35!black,title={Before you move on}]
What does \ml{തണുക്കുക} mean? (\textquotedblleft{}To grow cold; to cool down\textquotedblright{}.) Say it once more.
\end{tcolorbox}
\section[\texorpdfstring{\ml{ഇരമ്പുക} (irampuka)}{ഇരമ്പുക (irampuka)}]{\ml{ഇരമ്പുക} (irampuka) — to rumble; to roar}

\label{lesson:ML-C510-irampuka}



\begin{quote}
Before the new one: say the Malayalam for to cloud over, then the Malayalam for to grow cold.
\end{quote}

\subsection*{You'll want to know: \ml{ഇരമ്പുക}}
\textbf{\ml{ഇരമ്പുക}} — \emph{irampuka} — \textquotedblleft{}to rumble; to roar\textquotedblright{}.

\begin{quote}
\ml{രാത്രി} \ml{കടൽ} \ml{ഉറക്കെ} \ml{ഇരമ്പി}.
\end{quote}

\begin{quote}
\emph{rātri kaṭal uṟakke irampi.}
\end{quote}

\begin{quote}
At night the sea roared loudly.
\end{quote}

Past: \ml{ഇരമ്പി}, \emph{irampi}.

For the sea, the wind, engines and a crowd.

\begin{cousinweb}
\textbf{\ml{ഇരമ്പുക}} (\emph{irampuka}) is native Dravidian, \textquotedblleft{}to roar as the sea, to bluster\textquotedblright{} (DEDR 489).
\end{cousinweb}

\begin{grammarlens}[title={What you've built}]
One more word for getting around.
\end{grammarlens}

\subsection*{Your turn}
\begin{itemize}
\raggedright
  \item \emph{Say it:} \emph{irampuka}
  \item \emph{Say it:} \emph{irampuka}, once more
  \item \emph{Say it:} \emph{rātri kaṭal uṟakke irampi.} — At night the sea roared loudly.
  \item \emph{Recall:} say the Malayalam for to cloud over, then the Malayalam for to grow cold, then say \emph{irampuka} again
\end{itemize}

\begin{tcolorbox}[breakable,colback=teal!4,colframe=teal!35!black,title={Before you move on}]
What does \ml{ഇരമ്പുക} mean? (\textquotedblleft{}To rumble; to roar\textquotedblright{}.) Say it once more.
\end{tcolorbox}
